\documentclass{article}
\usepackage{graphicx} % Required for inserting images
\usepackage{amsfonts, amsmath, amssymb, amsthm, float}

\title{Meccanica Lagrangiana}
\author{Alessandro Ballini}
\date{June 2025}

\begin{document}

\maketitle

\tableofcontents
\newpage

\section{Moti unidimensionali}
I sistemi unidimensionali sono quelli descritti da una equazione differenziale ordinaria del secondo ordine della seguente forma
\begin{equation}
    m\ddot{x} = f(x, \dot{x}, t) \tag{1.1}
\end{equation}
con $f$ di classe $C^{\infty}$ e $m \in \mathbf{R}$. La $(1.1)$ può sempre essere riscritta come un sistema di due equazioni differenziali ordinarie del primo ordine in forma normale
\begin{equation}
    \begin{cases}
        \dot{x} = v \\
        m\dot{v} = f(x, \dot{x}, t)
    \end{cases} \tag{1.2}
\end{equation}
La teoria dei moti unidimensionali è l'unica che fornisce soluzioni generali. \\ \\
Tratteremo in particolare i sistemi \textit{autonomi}, ovvero sottoposti ad un campo di forze che non dipende esplicitamente dal tempo, e \textit{posizionali}, cioè il campo di forze non dipende dalla velocità del sistema, dunque $f = f(x), \ x \in \mathbf{R}$. \\
Definiamo l'energia totale del sistema $E$
\begin{equation}
    E = E(x, \dot{x}) = \frac{1}{2}m\dot{x}^2 + \phi(x) \tag{1.3}
\end{equation}
dove il primo termine del secondo membro è detta \textit{energia cinetica del sistema}, mentre $\phi : \mathbf{R} \longrightarrow \mathbf{R}$ è detto il \textit{potenziale} a cui è sottoposto il sistema, definito nel seguente modo
\begin{equation}
    \phi(x) = - \int f(x) \, dx \tag{1.4}
\end{equation}
Passiamo ora ad un risultato molto importante. \\ \\
\textbf{Teorema 1.1:} in un sistema descritto da un'equazione differenziale come $(1.1)$ l'energia totale è \textit{conservata}, ovvero
\begin{equation}
    \frac{d E}{d t} = 0 \tag{1.5}
\end{equation}

\textit{Dimostrazione} lungo le soluzioni del moto si ha
\begin{equation}
    \frac{d E}{d t} = m\dot{x}\ddot{x} + \frac{d\phi}{dx}\dot{x} = \dot{x}(m\ddot{x} - f(x)) = 0 \notag
\end{equation}
Possiamo dunque concludere, senza aver note le soluzioni, che l'energia totale è una costante del moto. \qed \\ \\
Partendo dalla definizione di energia si ha
\begin{equation}
    \frac{dx}{dt} = \pm \sqrt{\frac{2}{m}[E - \phi(x)]} \tag{1.6}
\end{equation}
Dalla teoria del calcolo differenziale sappiamo che $\frac{d}{dt}(f^{-1}) = \left( \frac{df}{dt} \right)^{-1}$. Inoltre, per intervalli di tempo abbastanza piccoli possiamo considerare $t \mapsto x(t)$ iniettiva e quindi invertibile. Pertanto vale la seguente relazione integrale
\begin{equation}
    t - t_0 = \pm \int_{x_0}^{x} \frac{dx'}{\sqrt{\frac{2}{m}[E - \phi(x')]}} \tag{1.7}
\end{equation}
$(1.7)$ permette di calcolare, a meno di integrare, il tempo impiegato per passare da una punto fissato $x_0$ ad un punto $x$, note le condizioni iniziali e il potenziale a cui è sottoposto il sistema. \\
Dato che $E$ è una costante del moto, è possibile determinarla evendo note solamente le condizioni iniziali $(x_0, \dot{x}_0)$, elemento dello spazio delle fasi. Dunque, ricordando $(1.3)$, si ha che una volta fissato il valore dell'energia totale $E$, il moto si svolge in una determinata regione dello spazio delle fasi. Infatti, il moto, intesto come il supporto $M$ di una curva nello spazio delle fasi, deve soddisfare la seguente proprietà
\begin{equation}
    M = \left\{ (x, \dot{x}) \in \mathbf{R}^2 : \frac{1}{2}m\dot{x}^2 + \phi(x) = E \right\} \tag{1.8}
\end{equation}
$(1.8)$ non è altro che una curva di livello per la funzione $E(x, \dot{x})$. Definiamo i punti di inversione $x_-, \ x_+$ dello spazio delle coordinate nel seguente modo
\begin{equation}
    \phi(x_-) = E, \quad \phi(x_+) = E \tag{1.9}
\end{equation}
ovvero i punti per i quali l'energia totale del sistema è dovuta interamente al potenziale, pertanto si ha $\dot{x}_-, \dot{x}_+ = 0$. \\
In genere un sistema può avere più punti di inversione, ciò dipende interamente dalla funzione $\phi$. Tuttavia, in tutti i casi, se la posizione iniziale $x_0 \in (x_-, x_+)$, tutto il moto è vincolato a restare entro i due punti di inversione, conseguenza del \textit{teorema 1.1}. \\ \\
Notiamo che per $x = x_+$ (o $x = x_-$) la funzione integranda in $(1.7)$ ha una singolarità in un estremo di integrazione, vediamo allora sotto quali condizioni possiamo affremare che l'integrale è finito. \\ \\
\textbf{Teorema 1.2:} sia $\phi'(x_+) = \alpha > 0$, allora
\begin{equation}
    \int_{x_0}^{x_+} \frac{dx}{\sqrt{\frac{2}{m}[E - \phi(x)]}} < \infty \tag{1.10}
\end{equation}

\textit{Dimostrazione} sia $\varepsilon > 0$ abbastanza piccolo, definiamo $\overline{x} = x_+ - \varepsilon$ e spezziamo l'integrale lungo i due intervalli $(x_0, \overline{x})$ e $(\overline{x}, x_+)$. L'integrazione lungo il primo intervallo non presenta problemi. Resta da verificare che
\begin{equation}
    \int_{\overline{x}}^{x_+} \frac{dx}{\sqrt{\frac{2}{m}[E - \phi(x)]}} < \infty \tag{1.11}
\end{equation}
Espandiamo il denominatore della funzione integranda al primo ordine con \textit{resto di Lagrange} attorno al punto di inversione $x_+$, ovvero
\begin{equation}
    E - \phi(x) = E - \phi(x_+) + [E - \phi(x)]'\left|_{x = s} \right.(x - x_+) \tag{1.12}
\end{equation}
per un ooportuno $s \in (\overline{x}, x_+)$. Notiamo subito che, per definizione di punto di inversione, si ha $E - \phi(x_+) = 0$, allora
\begin{equation}
    E - \phi(x) = -\phi'(s)(x - x_+) \tag{1.13}
\end{equation}
Sostituendo in $(1.11)$ si ottiene
\begin{equation}
    \int_{\overline{x}}^{x_+} \frac{dx}{\sqrt{\frac{2}{m}[E - \phi(x)]}} =  \int_{\overline{x}}^{x_+} \frac{dx}{\sqrt{-\frac{2}{m}\phi'(\xi)(x - x_+)}} < \infty \tag{1.14}
\end{equation}
per confronto. Analogamente si dimostra nel caso per $x_-$. \qed \\ \\
Vediamo ora sotto quali condizione si ha che $(1.7)$ diverge. \\ \\
\textbf{Teorema 1.3:} sia $\phi'(x_+) = 0$, allora
\begin{align}
    t - t_0 = \int_{x_0}^{x_+} \frac{dx}{\sqrt{\frac{2}{m}[E - \phi(x)]}} \tag{1.15}
\end{align}
diverge. \\

\textit{Dimostrazione:} seguendo un ragionamento identico a quello fatto per il \textit{teorema 1.2}, si ha che
\begin{equation}
    E - \phi(x) = E - \phi(x_+) + [E - \phi(x)]''\left|_{x = \xi}\right. (x - x_+)^2 \notag
\end{equation}
dato che $\phi'(x_+) = 0$. Segue che
\begin{equation}
    E - \phi(x) = -\phi''(s) (x - x_+)^2 \notag
\end{equation}
Tuttavia, sostituendo in $(1.7)$, si ottiene una funzione integranda proporzionale a $\frac{1}{x}$, quindi sappiamo che l'integrale diverge, per confronto. \qed \\ \\
Si vedrà che se $\phi'(x_+) = 0$ allora $x_+$ è un punto di equilibrio del sistema, ovvero è una soluzione costante, perciò, intuitivamente per l'unicità della soluzione dell'equazione differenziale, deve per forza valere il \textit{teorema 1.3}, infatti il sistema può solamente tendere asintoticamente a $x_0$, senza mai raggiungerlo.

\newpage
\section{Stabilità}
Siano $\Omega \in \mathbf{R}^n$, $f : \Omega \longrightarrow \mathbf{R}^n$ di classe $C^1$ (o superiore), $y_0 \in \Omega$ tale che $f(y_0) = 0$, allora $y_0$ è un \textit{punto di equilibrio} del sistema $\dot{y} = f(y)$. \\ \\
\textbf{Definizione:} sia 
\begin{equation}
    \begin{cases}
        \dot{y} = f(y) \\
        y(t_0) = y_0
    \end{cases} \notag
\end{equation}
un sistema autonomo con soluzione $\varphi( \cdot; t_0, y_0)$. Sia $r > 0$, allora $y_0$ è detto punto di \textit{equilibrio stabile} se: \\ \\
\textbf{1.} per ogni $p \in D_r(y_0)$ la soluzione $\varphi(\cdot; t_0, p)$ è definita almeno per ogni $t \in [0, +\infty)$. \\ \\
\textbf{2.} per ogni $\varepsilon > 0$ esiste $\delta_\varepsilon > 0$ tale che se $p \in D_{\delta_\varepsilon}(y_0)$ si ha
\begin{align}
    \|\varphi(t; t_0, p) - \varphi(t; t_0, y_0)\| < \varepsilon \tag{2.2}
\end{align}
per ogni $t \geq 0$. \\ \\
A parole, un punto di equilibrio $y_0$ è detto stabile se per ogni intorno arbitrariamente piccolo $U$ di $y_0$ è possibile trovare un punto $p$ in esso contenuto, tale che la soluzione associata a $p$ come condizione iniziale rimane all'interno di $U$ per ogni $t \geq t_0$. \\ \\ 
Un altro caso di stabilità sono i punti \textit{asintoticamente stabili}, vediamo come sono definiti. \\ \\
\textbf{Definizione:} un punto $y_0 \in \Omega$ è detto asintoticamente stabile se è stabile e se esiste $\varepsilon > 0$ tale che, per ogni $p \in D_\varepsilon(y_0)$, si ha
\begin{align}
    \lim_{t \rightarrow \infty} \|\varphi(t; t_0, p) - \varphi(t; t_0, y_0)\| = 0 \tag{2.3}
\end{align}
Ovvero un punto è asintoticamente stabile se, per punti vicini ad esso, il moto tende a quel punto, senza però mai raggiungerlo (per unicità della soluzione). \\ \\
Vediamo ora un importante risultato che permette di determinare la stabilità di un punto di equilibrio \\ \\
\textbf{Teorema (stabilità di Ljapunov):} sia $y_0$ un punto di equilibrio per $\dot{y} = f(y)$, sia poi $W : \mathbf{R}^n \longrightarrow \mathbf{R}$ continua tale che $y_0$ è un punto di minimo isolato per $W$ e $\dot{W} \leq 0$ lungo le soluzioni del moto. Allora $y_0$ è un punto di equilibrio stabile. \\

\textbf{Dimostrazione:} poiché ci interessa solamente che $y_0$ sia un minimo per $W$ possiamo, a meno di ridefinire la funzione, scegliere $W$ tale che $W(y_0) = 0$ e $W(y) > 0$ per ogni $y \in \mathbf{R}^n \setminus \{y_0\}$. Lungo le soluzioni del moto si ha
\begin{equation}
    \frac{dW}{dt} = \nabla W \cdot \dot{y} = \nabla W \cdot f(y) \leq 0 \tag{2.4}
\end{equation}
Siano $c > 0$ e $\mathcal{G}_c = \{ y \in \mathbf{R}^n : W(y) \leq c \}$. Dato che $W$ è continua si ha che $\mathcal{G}_c$ è compatto. Pertanto, l'ipotesi $\dot{W} \leq 0$ lungo le soluzioni del moto implica, per continuità della funzione $W$, che anche la soluzione resti arbitrariamente vicina a $y_0$. Infatti sia $\varphi(\cdot; t_0, y)$ una soluzione del sistema con $y \in \mathcal{G}_c$, allora $W(\varphi(t)) \leq c$ per ogni $t \geq t_0$, quindi, per ipotesi di continuità, anche $\varphi$ resta arbitrariamente vicina al punto $y_0$ per ogni $t \geq t_0$. \\ \\
\textbf{Teorema (di Dirichlet):} sia $x_0$ un minimo locale per il potenziale $\phi$, allora $x_0$ è un punto di equilibrio stabile per il sistema $\ddot{x} = -\nabla \phi$. \\

\textbf{Dimostrazione:} dato che $x_0$ è un punto di minimo per il potenziale si ha che $\nabla \phi (x_0) = 0$, dunque in $x_0$ il campo di forze è nullo. Inoltre, $x_0$ è un minimo isolato dell'energia, posto $\dot{x}(t_0) = 0$. Quindi $W = E$ è una funzione di Ljapunov. \\ \\
Studiamo ora come evolve il sistema per piccole oscillazioni attorno ad un punto di equilibrio stabile. Vediamo come si comporta il potenziale attorno ad un punto stabile. \\ \\
\textbf{Proposizione 2.1:} sia $\phi : \mathbf{R}^n \longrightarrow \mathbf{R}$ il potenziale di un sistema autonomo, sia $x_0$ un punto di equilibrio stabile per il sistema tale che $\phi''(x_0) = \alpha > 0$. Allora per $x \rightarrow x_0$ si ha che $\phi(x) \sim \frac{1}{2}\alpha(x - x_0)^2$. \\

\textbf{Dimostrazione:} poniamo $\phi(x_0) = 0$, poiché $x_0$ è un punto di equilibrio si ha che $\phi'(x_0) = 0$, quindi
\begin{equation}
    \phi(x) = \frac{1}{2}\alpha(x - x_0)^2, \quad x \longrightarrow x_0 \tag{2.5}
\end{equation}
ovvero il potenziale nelle vicinanze del punto di equilibrio stabile si comporta come il potenziale dell'oscillatore armonico.

\newpage
\section{Formalismo Lagrangiano}
Consideriamo ora sistemi discreti a più gradi di libertà, quindi descritti da un sistema di equazioni differenziali del tipo
\begin{equation}
    m\ddot{\mathbf{x}} = - \nabla \phi \tag{3.1}
\end{equation}
dove $\mathbf{x} : \mathbf{I} \subset \mathbf{R} \longrightarrow \mathbf{R}^n$, cammino liscio, e $\phi : \Omega \subset \mathbf{R}^n \longrightarrow \mathbf{R}^n$ potenziale di classe $C^2$. \\ \\
Si mostra che le soluzioni del moto sono quelle che minimizzano una quantità detta \textit{azione} $S$, almeno in prima approssimazione. L'azione è definita nel seguente modo
\begin{equation}
    S = \int_{t_1}^{t_2} L(\mathbf{x}(t), \dot{\mathbf{x}}(t), t) \, dt \tag{3.2}
\end{equation}
Sia $\boldsymbol{\delta}$ una variazione del cammino ad estremi fissati, ovvero $\boldsymbol{\delta}(t_1) = 0 = \boldsymbol{\delta}(t_2)$. Vediamo che condizioni deve soddisfare $L$ affinché l'azione sia minimizzata.
\begin{equation}
    \delta S = \int_{t_1}^{t_2} \left( L(\mathbf{x} + \boldsymbol{\delta}, \dot{\mathbf{x}} + \dot{\boldsymbol{\delta}}, t) - L(\mathbf{x}, \dot{\mathbf{x}}, t) \right) \, dt = 0 \tag{3.3}
\end{equation}
equivale a
\begin{equation}
    \delta S = \int_{t_1}^{t_2} \left(\nabla_{\mathbf{x}}L \cdot \boldsymbol{\delta} + \nabla_{\dot{\mathbf{x}}} \cdot \dot{\boldsymbol{\delta}}\right) \, dt = 0 \tag{3.3}
\end{equation}
In particolare $\frac{d}{dt}\delta = \dot{\delta}$, quindi, integrando per parti il secondo termine, si ha
\begin{equation}
    \int_{t_1}^{t_2} \nabla_{\mathbf{x}}L \cdot \boldsymbol{\delta} \, dt = \int_{t_1}^{t_2} \nabla_{\dot{\mathbf{x}}}L \cdot \boldsymbol{\delta} \, dt - \nabla_{\dot{\mathbf{x}}}L \cdot \boldsymbol{\delta}|_{t_1}^{t_2} \tag{3.4}
\end{equation}
Il secondo termine del secondo membro è nullo, infatti la variazione è ad estremi fissati. Dato che $(3.4)$ deve valere per ogni variazione $\boldsymbol{\delta}$ arbitraria, la funzione $L$ deve soddisfare il seguente sistema di equazioni differenziali, dette \textit{equazioni di Eulero - Lagrange}
\begin{equation}
    \frac{\partial L}{\partial x_i} - \frac{d}{dt}\frac{\partial L}{\partial \dot{x}_i} = 0 \qquad i = 1, 2, 3, ... , n \tag{3.5}
\end{equation}
Esse costituiscono un sistema di $n$ equazioni differenziali del secondo ordine, dipendenti da $2n$ costanti che sono determinabili a partire dalle condizioni iniziali. \\ \\
La principale comodità del formalismo Lagrangiano sta nella sua dipendenza da una condizione di minimizzazione. Qualsiasi sia il set di variabili dinamiche indipendenti $\mathbf{q} = (q_1, \dots, q_n)$ (legate alle vecchie coordinate da un diffeomorfismo) scelte per descrivere il sistema, l'azione $S$ sarà ancora minimizzata. Dunque, è possibile scrivere le equazioni di Eulero-Lagrange nella seguente forma
\begin{equation}
    \frac{\partial L}{\partial q_i} - \frac{d}{dt}\frac{\partial L}{\partial \dot{q}_i} = 0 \qquad i = 1, 2, 3, ... , n \tag{3.6}
\end{equation}
Per sistemi meccanici è facile dimostrare che la funzione $L$ prende la seguente forma
\begin{equation}
    L(\mathbf{q}, \dot{\mathbf{q}}, t) = T(\mathbf{q}, \dot{\mathbf{q}}) - \phi(\mathbf{q}) \tag{3.7}
\end{equation}
Da $(4.1)$ risulta chiaro come l'utilizzo di un set di coordinate piuttosto che un altro non influenza la variazione prima del funzionale azione. \\ \\
\textbf{Esempio:} in un sistema bidimensionale formato da una particella in moto libero, passando da coordinate cartesiane a coordinate polari ponendo
\begin{equation}
    \begin{cases}
        x = \rho \cos \theta \\
        y = \rho \sin \theta
    \end{cases} \notag
\end{equation}
l'energia cinetica $T$ diventa
\begin{equation}
    T = \frac{1}{2}m\rho^2 \dot{\theta}^2 \tag{3.8}
\end{equation}
in tal caso la variazione prima dell'azione è ancora nulla e quindi le equazioni del moto restano invariate.

\newpage
\section{Sistema linearizzato}
Come visto nel capitolo $2$, attorno a punti di equilibrio stabile il sistema si comporta come un oscillatore armonico, dunque viene comodo riscrivere le equazioni del moto in tale condizioni per studiare il sistema. \\
Definiamo una nuova variabile meccanica che tiene conto dello spostamento relativo alla posizione di equilibrio. Sia $x_0$ un punto di equilibrio stabile, poniamo
\begin{equation}
    \eta = x - x_0 \qquad \dot{\eta} = \dot{x} \tag{5.1}
\end{equation}
Riscriviamo la funzione di Lagrange nel seguente modo
\begin{equation}
    L(\eta, \dot{\eta}) = \frac{1}{2}\sum_{i,j} T_{i,j}(\eta)\dot{\eta}_i \dot{\eta}_j - \phi(\eta) \tag{5.2}
\end{equation}
Selezioniamo ora un intorno abbastanza piccolo del punto di equilibrio nello spazio delle fasi tale che è possibile trascurare le correzioni di ordine maggiore di $2$, possiamo allora approssimare $L$ con
\begin{equation}
    L = \frac{1}{2}\sum_{i,j}T_{i,j}(0)\dot{\eta}_i \dot{\eta}_j - \frac{1}{2}\sum_{i,j}\phi_{i,j}(0) \eta_i \eta_j \tag{5.3}
\end{equation}
dove si è posto
\begin{equation}
    \phi_{i,j} = \frac{\partial^2 \phi}{\partial \eta_i \partial \eta_j} \tag{5.4}
\end{equation}
Supponendo sufficiente regolarità, per il \textit{teorema di Schwartz}, si ha $\phi_{i,j}$ è una matrice simmetrica. Inoltre, dato che $\eta = 0$ è un punto di equilibrio stabile, quindi di minimo isolato per il potenziale, segue che $\phi_{i,j}$ è definita positiva. La matrice dell'energia cinetica $T_{i,j}$ anche è definita positiva, infatti $T \geq 0$ sempre, e simmetrica. \\
Le equazioni di Eulero-Lagrange associate sono
\begin{equation}
    T_{i,j}(0)\ddot{\eta} + \phi_{i,j}(0)\eta = 0 \tag{5.5}
\end{equation}
Notiamo subito che il sistema è linearizzato. Inoltre, per quanto visto nel capitolo $2$, sappiamo che il sistema si comporta come un oscillatore armonico attorno a $\eta = 0$, perciò è ragionevole cercare una soluzione oscillante della forma
\begin{equation}
    q(t) = \vec{u}\cos \omega t \tag{5.6}
\end{equation}
con $\vec{u} \in \mathbf{R}^n$ e $\omega \in \mathbf{R}$. Sostituendo nel sistema linearizzato si ha
\begin{equation}
    \left( \phi_{i,j} - \omega^2T_{i,j} \right) q(t) = 0 \quad \forall t \geq t_0 \tag{5.7}
\end{equation}
Abbiamo dunque ottenuto un'equazione differenziale agli autovalori, pertanto è possibile determinare la pulsazione $\omega$ del sistema dal polinomio caratteristico dato dalla condizione $\det\left( \phi_{i,j} - \omega^2T_{i,j} \right) = 0$.

\newpage
\section{Il teorema di E. Noether}
Sia $X$ uno spazio topologico, una simmetria è una biezione $f : X \longrightarrow X$. A parole, una simmetria è una mappa che lascia invariato l'oggetto sul quale viene applicata. \\
Una simmetria banale è l'identità $\mathbf{1}_X : X \longrightarrow X$. Un altro esempio è un qualsiasi tipo di rotazione attorno ad un asse passante per il centro di una sfera. \\ \\
Supponiamo di avere un sistema meccanico, descritto da una funzione di Lagrange $L$, \textit{invariante per traslazioni temporali}, cioè la biezione $t \mapsto t + \varepsilon$, dove $\varepsilon > 0$. \\
Allora si ha che la funzione di Lagrange $L$ non dipende esplicitamente dal tempo, ovvero $\frac{\partial L}{\partial t} = 0$, dunque lungo le soluzioni del moto si ottiene
\begin{equation}
    \frac{dL}{dt} = \nabla_{\mathbf{q}}L \cdot \dot{\mathbf{q}} + \nabla_{\dot{\mathbf{q}}}L \cdot \ddot{\mathbf{q}} = \frac{d}{dt}\nabla_{\dot{\mathbf{q}}}L \cdot \dot{\mathbf{q}} + \nabla_{\dot{\mathbf{q}}}L \cdot \ddot{\mathbf{q}} = \frac{d}{dt}\left(\nabla_{\dot{\mathbf{q}}}L \cdot \dot{\mathbf{q}}\right) \tag{6.1}
\end{equation}
$(6.1)$ può essere riscritto nel seguente modo
\begin{equation}
    \frac{d}{dt}\left(\nabla_{\dot{\mathbf{q}}}L \cdot \dot{\mathbf{q}} - L\right) = 0 \tag{6.2}
\end{equation}
ovvero
\begin{align}
    \frac{d}{dt}E = 0 \tag{6.3}
\end{align}
Concludiamo che per sistemi invarianti rispetto a traslazioni temporali l'energia è conservata, ciò sussiste anche per traslazioni temporali non infinitesime. \\ \\
Supponiamo di avere ora un sistema meccanico \textit{invariante per traslazioni spaziali}, ovvero $\mathbf{q} \mapsto \mathbf{q} + \delta\mathbf{q}$. Si ha che
\begin{equation}
    L' = L + \nabla_{\mathbf{q}}L \cdot \delta\mathbf{q} \tag{6.4}
\end{equation}
verifica ancora il principio di minima azione. Lungo le soluzioni del moto si ha che
\begin{equation}
    \nabla_{\mathbf{q}}L = \frac{d}{dt}\nabla_{\dot{\mathbf{q}}}L \tag{6.5}
\end{equation}
Sostituendo $(6.5)$ in $(6.4)$ si ottiene
\begin{equation}
    \frac{d}{dt}\nabla_{\dot{\mathbf{q}}}L \cdot \delta\mathbf{q} = 0 \tag{6.6}
\end{equation}
per ogni $\delta\mathbf{q}$ (arbitrario), ovvero
\begin{align}
    \mathbf{p} = \nabla_{\dot{\mathbf{q}}}L = \text{costante} \tag{6.7}
\end{align}
Pertanto, in sistemi invarianti per traslazioni spaziali rigide si conserva la quantità di moto. Ciò vale anche per singole componenti, infatti sarebbe corretto scrivere
\begin{equation}
    \frac{d}{dt}\frac{\partial L}{\partial \dot{q}_i} \delta q_i = 0 \quad \forall \delta q_i \implies p_i = \frac{\partial L}{\partial \dot{q}_i} = \text{costante} \notag
\end{equation}
Quindi possono essere costanti anche solamente alcune delle componenti della quantità di moto. \\ \\
Abbiamo notato con i precedenti esempi che se il sistema ammette una simmetria esiste una quantità conservata. Vediamo ora la formulazione generale di tale risultato. \\ \\
\textbf{Definizione:} una trasformazione di coordinate $\mathbf{Q} = f(\mathbf{q})$ è detta \textit{ammissibile} per il sistema se e solo se la funzione di Lagrange è invariante per tale trasformazione, cioè
\begin{equation}
    L(\mathbf{Q}, \dot{\mathbf{Q}}) = L(\mathbf{q}, \dot{\mathbf{q}}) \tag{6.8}
\end{equation}
\textbf{Definizione:} un \textit{gruppo ad una parametro} $s \in \mathbf{R}$ è una famiglia di trasformazioni invertibili
\begin{equation}
    H_s = \{ h_s : \mathbf{R}^n \longrightarrow \mathbf{R}^n : Q = h_s(q),\ s \in \mathbf{R} \} \tag{4.3.2}
\end{equation}
che soddisfano le seguenti proprietà
\begin{align}
    h_0(\mathbf{q}) &= \mathbf{q} \tag{6.9} \\
    (h_{s_1} \circ h_{s_2})(\mathbf{q}) &= h_{s_1 + s_2}(\mathbf{q}) \tag{6.10}
\end{align}
per ogni $\mathbf{q} \in \mathbf{R}^n$ e ogni $s_1, s_2 \in \mathbf{R}$. \\ \\
Se per ogni $s \in \mathbf{R}$ la trasformazione $h_s$ è ammissibile per il sistema, si dice che il gruppo $H_s$ è ammissibile. \\ \\
\textbf{Osservazione:} $(6.9)$ e $(6.10)$ comportano
\begin{equation}
    Q = h_s(q) \implies q = h_{-s}(Q) \tag{6.11}
\end{equation}
Vediamo ora il risultato principale di questo capitolo. \\ \\
\textbf{Teorema (di E. Noether):} se un sistema lagrangiano $L$ ammette un gruppo ad un parametro $H_s$, le equazioni di EL associate a $L$ ammettono un integrale primo $I$ dato da
\begin{equation}
    I(\mathbf{q}, \dot{\mathbf{q}}) = \sum_i \frac{\partial L}{\partial \dot{q}_i}\frac{\partial h_s^i}{\partial s} \tag{6.12}
\end{equation}

\textbf{Dimostrazione:} sia $\mathbf{q}(t)$ una soluzione del moto, per ipotesi di invarianza di $L$ rispetto al gruppo di trasformazioni si ha che $\mathbf{Q}(t, \sigma) = h_\sigma(\mathbf{q})$ è ancora una soluzione del moto per ogni $\sigma \in \mathbf{R}$. Segue che
\begin{equation}
    \frac{d}{dt}\nabla_{\dot{\mathbf{Q}}} L(\mathbf{Q}, \dot{\mathbf{Q}}) = \nabla_\mathbf{Q} L(\mathbf{Q}, \dot{\mathbf{Q}}) \tag{6.13}
\end{equation}
Ovvero $L(\mathbf{Q}, \dot{\mathbf{Q}})$ verifica ancora le equazioni di EL. Inoltre, l'ipotesi di invarianza comporta la seguente condizione
\begin{equation}
    0 = \frac{\partial}{\partial \sigma}L(\mathbf{Q}, \dot{\mathbf{Q}}) = \nabla_\mathbf{Q} L \cdot \frac{\partial \mathbf{Q}}{\partial \sigma} + \nabla_{\dot{\mathbf{Q}}} L \cdot \frac{\partial \dot{\mathbf{Q}}}{\partial \sigma} \tag{6.14}
\end{equation}
Sostituendo $(6.13)$ in $(6.14)$ e moltiplicando entrambi i membri per $\frac{\partial\mathbf{Q}}{\partial\sigma}$ si ottiene
\begin{equation}
    \frac{d}{dt}\nabla_{\dot{\mathbf{Q}}} L \cdot \frac{\partial \mathbf{Q}}{\partial \sigma} + \nabla_{\dot{\mathbf{Q}}}L \cdot \frac{\partial \dot{\mathbf{Q}}}{\partial \sigma} = 0 \tag{6.15}
\end{equation}
Pertanto, siccome supponiamo abbastanza regolare la trasformazione $h_s$ e quindi è possibile invertire l'ordine di derivazione, possiamo scrivere $(6.15)$ come segue
\begin{equation}
    \frac{d}{dt}\left( \nabla_{\dot{Q}}L \cdot \frac{\partial Q}{\partial \sigma} \right) = 0 \tag{6.16}
\end{equation}
che, a meno di sviluppare il prodotto scalare, calcolato per $\sigma = 0$, non è altro che la tesi. \qed
\end{document}